% ============================================================================
%  Rozdział 6. Zakończenie i wnioski
% ============================================================================
\chapter{Zakończenie i wnioski}
\label{rozdz:zakonczenie}

W~ramach pracy zaprojektowano i~zaimplementowano trójwymiarową grę zręcznościową
w~silniku Unity. Gra zawiera trzy poziomy, kompletny przepływ od menu
głównego przez ekran ładowania do rozgrywki i~okna końca poziomu, system punktów życia,
dwa typy przeciwników, sześć mechanizmów zmieniających ruch kuli (pola przyspieszające,
odbijacze, obracające się belki, ruchome ściany, wiatr i~przycisk \emph{Push}), minimapę
oraz pomiar czasu z~zapisem statystyk. Gra działa w~edytorze Unity oraz na urządzeniach
z~systemem Android w~wersji 9 lub nowszej.

\section{Podsumowanie realizacji celów}

Cel główny pracy został osiągnięty. Zrealizowano mechanikę zbierania monet
z~warunkiem odblokowania mety, sterowanie gestem przeciągnięcia, odbicia i~tarcie, system
przeciwników o~zachowaniach autonomicznych (patrol kaktusów) oraz punkty życia z~leczeniem.
Architekturę oparto na wstrzykiwaniu zależności z~biblioteką Zenject i~podziale kodu na
50~zestawów kompilacji, a~logikę rozgrywki oddzielono od komponentów silnika. Interfejs
użytkownika zbudowano w~podejściu MVVM, a~komunikację pomiędzy modułami oparto na
strumieniach biblioteki R3 i~operacjach asynchronicznych UniTask. Efekty cząsteczkowe
korzystają z~puli obiektów, a~monety z~analogicznego mechanizmu wielokrotnego użycia.

Część założeń wniosku zrealizowano w~zmienionej formie. Zamiast biblioteki UniRx
wykorzystano R3, którą autor UniRx wskazuje jako jej rozwiniętą wersję; zmiana ta nie
wpłynęła na architekturę, ponieważ obie biblioteki opierają się na tych samych pojęciach.
Platformą docelową stał się Android, ponieważ sterowanie gestem okazało się naturalne dla
ekranu dotykowego, przy czym gra pozostaje w~pełni grywalna w~edytorze przy użyciu myszy.
Zapis danych obejmuje lokalne statystyki poziomów, a~integracja z~Firebase została
ograniczona do inicjalizacji i~weryfikacji zależności. Przygotowana obsługa reklam
z~nagrodą Appodeal nie została jeszcze podłączona do rozgrywki.

\section{Najważniejsze wnioski}

Pierwszy wniosek dotyczy kosztów i~korzyści wstrzykiwania zależności w~grze tej skali.
Kod gry liczy 256~plików, z~czego znaczną część stanowią bindery, klasy danych i~komponenty
znacznikowe. Dla małej gry jest to narzut wyraźny, jednak w~zamian każda klasa jawnie
deklaruje swoje zależności, a~żadna nie wyszukuje obiektów w~scenie. Dodanie w~późnej fazie
projektu wiatru, minimapy, licznika czasu, apteczek i~trzech typów przeszkód wymagało
jedynie niewielkich rozszerzeń istniejących klas, takich jak operacje \kod{Launch}
i~\kod{Knock} oraz prędkość dryfu w~logice ruchu kuli czy metoda leczenia w~klasie
\klasa{HealthStorage}, a~nowe moduły zostały podłączone przez kolejne bindery
w~istniejących instalatorach. Krótkie klasy o~jednej
odpowiedzialności, zgodne z~zaleceniami McConnella \cite{mcconnell}, ułatwiły też
lokalizowanie błędów zgłaszanych w~sesjach testowych.

Drugi wniosek dotyczy programowania reaktywnego. Strumienie R3 pozwoliły zapisać łańcuchy
zdarzeń, takie jak droga od zebrania monety do okna końca poziomu, w~sposób, w~którym
nadawca nie zna odbiorców. Wymaga to jednak dyscypliny: każda subskrypcja musi zostać
zwolniona, a~gdy wynik zależy od kolejności reakcji kilku odbiorców, kolejność subskrypcji
staje się częścią kontraktu, co w~projekcie trzeba było udokumentować w~komentarzu.

Trzeci wniosek dotyczy fizyki. Własny model ruchu oparty na kierunku i~mocy, nałożony na
symulację PhysX, dał przewidywalne odbicia i~drogę hamowania, którą można obliczyć
z~parametrów konfiguracji. Przeniesienie wszystkich wartości liczbowych do plików JSON
okazało się szczególnie przydatne, ponieważ większość poprawek po sesjach testowych
dotyczyła strojenia parametrów, a~nie zmian w~kodzie.

Czwarty wniosek wynika z~testów. Automatyczny scenariusz uruchomieniowy ujawnił
ostrzeżenie Zenject związane z~tworzeniem kuli w~trakcie instalacji powiązań, a~analiza
modelu sterowania pokazała, że moc strzału zależy od odległości mierzonej w~pikselach.
Na ekranach o~różnej gęstości pikseli ten sam fizyczny gest daje więc różną moc. Oba
problemy nie utrudniają gry, ale wskazują miejsca, które należy poprawić przed dalszym
rozwojem.

\section{Dalsze prace i możliwości rozwoju}

Najbliższe prace powinny objąć usprawnienia wskazane w~testach: normalizację długości
gestu względem rozmiaru lub gęstości pikseli ekranu, przeniesienie tworzenia kuli do etapu
inicjalizacji kontenera, obsługę wyczerpania puli efektów oraz eliminację alokacji przy
aktualizacji licznika czasu. Klasy logiki nie zależą od komponentów silnika, więc można je
objąć testami jednostkowymi w~trybie edycji pakietu Unity Test Framework bez uruchamiania
scen. Warto również wykonać systematyczne pomiary wydajności na kilku telefonach.

Z~punktu widzenia gracza naturalnymi kierunkami rozwoju są: warstwa dźwiękowa, lokalizacja
interfejsu na język polski, kolejne poziomy oraz edytor poziomów wykorzystujący istniejący
mechanizm punktów startowych. Firebase, już zainicjalizowany w~grze, umożliwia dodanie
zapisu postępów w~chmurze i~globalnych tablic najlepszych czasów, a~przygotowana klasa
reklam z~nagrodą może posłużyć do przywracania zdrowia w~zamian za obejrzenie reklamy.

\section{Obszary zastosowania opracowanych rozwiązań}

Zastosowana architektura może posłużyć jako szablon dla innych niewielkich gier mobilnych
tworzonych w~Unity. Elementy takie jak wczytywanie konfiguracji z~walidacją przed
budowaniem, binder jako jednostka rejestracji powiązań, generyczne klasy MVVM dla prostych
paneli interfejsu czy zautomatyzowane budowanie w~osobnej kopii projektu nie zależą od
mechaniki tej konkretnej gry. Projekt może być także materiałem dydaktycznym pokazującym
zastosowanie wstrzykiwania zależności i~programowania reaktywnego w~silniku, w~którym
podejścia te nie są narzucane przez sam silnik.
